\section{Related work}\label{sec:related}

\subsection{Probabilistic logic programming}

ProbLog~\cite{deraedt2007} attaches probabilities to logic programs
under the \emph{distribution semantics}: a program denotes a measure over
possible worlds, inference computes exact marginals via knowledge
compilation to weighted Boolean formulas~\cite{fierens2015}, and the
system supports conditioning on evidence --- an operation \pla{}
deliberately lacks. (Both capabilities are verified against the
installed \texttt{problog} package by our test suite.) aProbLog
generalizes the semantics to arbitrary commutative
semirings~\cite{kimmig2011}, which is the correct formal umbrella for
non-probabilistic confidence calculi such as \pla{}'s. DeepProbLog adds
neural predicates with end-to-end gradient training~\cite{manhaeve2018},
and Scallop compiles differentiable reasoning over provenance
semirings~\cite{li2023}. All of these systems ship as Python packages;
\pla{} claims no language-based novelty whatsoever.

\subsection{Weighted logic and statistical relational learning}

Markov Logic Networks attach real-valued weights to first-order clauses
defining a Markov random field~\cite{richardson2006}; Probabilistic Soft
Logic relaxes truth values to $[0,1]$ with hinge-loss MRFs, making it
the nearest formally grounded cousin of \pla{}'s soft
confidences~\cite{bach2017}; Logical Neural Networks make weighted
real-valued logic differentiable while maintaining truth-value
bounds~\cite{riegel2020}; recent surveys map this converging
territory~\cite{marra2024,garcez2023}. In all of these, rule weights are
\emph{static at inference time}: adapting a rule's reliability to a
changing operational context is retraining's job, not the semantics'.
That gap --- context-conditioned weights inside an interpretable
calculus --- is the opening \pla{} targets, and it connects to the
concept-drift literature~\cite{gama2014}, where adaptation is studied
outside logical representations. Bayesian networks (e.g.\ the
educational-standard pgmpy package~\cite{ankan2015}) provide
joint-distribution posteriors --- verified against the installed package
--- but not rule-shaped, trace-first explanations.

A related, newer line uses crisp symbolic solvers to verify the
reasoning of large language models~\cite{pan2023,olausson2023,ye2023},
following the generate-and-verify architecture argued for
in~\cite{kambhampati2024}. A verifier that propagates \emph{uncertain}
rule confidences rather than crisp truth is an open slot in that
architecture which \pla{}-style calculi could fill; we leave this as
future work.

\subsection{Uncertainty calculi: \pla{}'s lineage}

\pla{}'s operators descend from MYCIN's certainty
factors~\cite{shortliffe1975}: rule-attached confidences, attenuation by
the weakest antecedent --- Zadeh's G\"odel t-norm~\cite{zadeh1965} ---
and parallel combination of co-supporting rules. Indeed \pla{}'s default
noisy-OR \emph{is} the CF parallel-combination formula
$\mathit{CF} = \mathit{CF}_1 + \mathit{CF}_2(1-\mathit{CF}_1)$, an
equivalence enforced by a unit test. Heckerman showed that the CF
calculus is probabilistically coherent only under restrictive
independence and modularity assumptions~\cite{heckerman1986}, and the
field's retrospective explains the subsequent migration to belief
networks~\cite{heckerman1992}. \pla{}'s response is not to dispute the
critique but to make its subject explicit: it drops the probability
claim entirely (Section~\ref{sec:notprob}), names the independence
stance as a per-knowledge-base operator choice, and provides log-odds
modes that match the likelihood-ratio special case Heckerman identified
as coherent.

\subsection{Positioning}

\pla{} does not compete with the systems above on semantic guarantees;
it occupies the point they leave open: a fully readable engine whose
primary output is the reasoning trace, whose rule weights are
context-conditioned and learnable, and whose explanation fidelity is
measured rather than asserted~\cite{rudin2019,jacovi2020}. The intended
domains are education and audit-style settings~\cite{appelbaum2017},
where inherent interpretability is a requirement, not a preference, and
where fraud-detection practice~\cite{west2016,dalpozzolo2015,bao2020}
supplies public data for evaluation.
